\subsection{An unconditional bound for the actual cubic average polynomial}

\begin{proposition}[P-F42: a Weyl--cone bound on the actual coefficient norm]
\label{prop:cubicaverageupper}
Use the ordinary cubic stretching function $f_n$, average polynomial
$B_n$, and coefficient norm $U_n$ of P-F34. For every $n\ge4$, put
\[
 m=\binom{n+2}{3},\qquad a=m-n.
\]
Then
\[
 U_n\le
 8n!\,m^{3n}\bigl(288n^3\,3^n\bigr)^{m-n}.
 \tag{P-F42a}
\]
In particular,
\[
 \log_2 U_n\le
 \frac{\log_2 3}{6}\,n^4+O(n^3\log n),\qquad
 \limsup_{n\longrightarrow\infty}
       \frac{\log_2 U_n}{n^4}
       \le\frac{\log_2 3}{6}.
 \tag{P-F42b}
\]
The implied constant is absolute. This is an unconditional estimate
for the actual cubic invariant ring, including its Weyl alternation.
Its leading constant is greater than $1/24$, so it does not establish
the sufficient remaining condition in P-F34 or refute PH2 or SH.
It is a coefficient bound, not a claim of a polynomial-time procedure
for computing $B_n$.
\end{proposition}

\begin{proof}
\emph{1. The exact vector-partition formula.}
Let
\[
 g=\gcd(n,3),\quad c=3/g,\quad d_0=n/g,\quad
 \delta=(n-1,n-2,\ldots,0).
\]
Index the columns of an $n$ by $m$ matrix $A$ by all degree-three
monomials in $n$ variables, and let a column $\alpha_i$ be that
monomial's exponent vector. Each column is nonnegative with
coordinate sum three and Euclidean norm at most three.
The columns $3e_j$ show that $A$ has rank $n$. Define
\[
 p(b)=\#\{\beta\in\mathbb Z_{\ge0}^{m}:A\beta=b\}.
\]
Then
\[
 f_n(s)=
 \sum_{\sigma\in S_n}\operatorname{sgn}(\sigma)\,
       p(cs\mathbf1+b_\sigma),\qquad
 b_\sigma=\delta-\sigma\delta.
 \tag{P-F42c}
\]
In particular $|(b_\sigma)_j|\le n-1$ and
$\sum_j(b_\sigma)_j=0$.

Here is the coefficient extraction proving (P-F42c).
The monomial character of
$\operatorname{Sym}^{d_0s}(\operatorname{Sym}^3\mathbb C^n)$
counts the vectors $\beta$ with $\sum_i\beta_i=d_0s$.
If $A\beta=cs\mathbf1+b_\sigma$, the last equation is automatic:
summing coordinates gives $3\sum_i\beta_i=ncs=3d_0s$.
Multiply this character by the alternant
$\sum_{\sigma}\operatorname{sgn}(\sigma)x^{\sigma\delta}$
and extract the coefficient of $x^{cs\mathbf1+\delta}$.
The Weyl character formula says that an irreducible summand with
highest weight $\lambda$ contributes its alternant with exponents
the permutations of $\lambda+\delta$.
The vector $cs\mathbf1+\delta$ is strictly decreasing, so it occurs
only when $\lambda=cs\mathbf1$, and then with coefficient one.
This is precisely the $SL_n$-invariant summand.
The Weyl formula used here is also displayed in GCT VI \cite{PGCTVI},
Section~5.1, equation~(5.1). No cancellation or estimate has yet
been discarded.

\emph{2. One perturbation works for every integer right-hand side.}
We shall bound the average constituent polynomial of each function
$p(cs\mathbf1+b_\sigma)$ separately. Choose a rational vector
$\epsilon\in(0,1)^m$ as follows.
For each rational hyperplane spanned by a rank-$(n-1)$ subset
of columns of $A$, choose a nonzero primitive integer normal
$\ell_H$. There are finitely many such hyperplanes. Take
$\epsilon$ sufficiently small and outside their inverse-image
hyperplanes so that
\[
             0<|\ell_H A\epsilon|<1
                    \quad\hbox{for every }H.
 \tag{P-F42d}
\]
This is possible: $A$ has full rank, so each linear functional
$\ell_H A$ is nonzero, and a finite union of proper hyperplanes
cannot cover an open positive box. Rational points outside that
union exist in a sufficiently small such box.
For every integer vector $b$, the vector $b+A\epsilon$ lies on
none of the specified hyperplanes, since $\ell_H b$ is an integer
and the nonzero number $\ell_H A\epsilon$ has absolute value less
than one.

For $b=cs\mathbf1+b_\sigma$, consider
\[
 P_s^\epsilon=
 \{\beta\in\mathbb R^m:A\beta=b,\ \beta_i\ge-\epsilon_i\}.
 \tag{P-F42e}
\]
Its integer points are exactly the nonnegative integer solutions
of $A\beta=b$. It is bounded, because the column sums give
$\sum_i\beta_i=d_0s$ and every coordinate has a lower bound.
For all sufficiently large $s$, it is nonempty of dimension
$a=m-n$: the translated right-hand side $b+A\epsilon$ has
strictly positive coordinates, and the columns $3e_j$ allow a
solution with every translated variable $u_i=\beta_i+\epsilon_i$
strictly positive.
At a vertex, the positive coordinates of $u$ must correspond
to linearly independent columns. There must be exactly $n$ of
them, since otherwise $b+A\epsilon$ would lie in one of the
hyperplanes excluded above. Thus the polytope is simple, with
vertices indexed by feasible $n$-column bases.
For a fixed basis, every basic coordinate is affine in $s$.
Its strict positivity is therefore eventually constant.
There is consequently one threshold beyond which the set of
feasible bases is constant, simultaneously for all the finitely
many shifts $b_\sigma$. No bound on that threshold is needed.

\emph{3. Bounds on the bases and one common generic direction.}
For an $n$-column basis $J$, write
\[
 C=A_J,\qquad d=|\det C|,\qquad I=\{1,\ldots,m\}\setminus J.
\]
Hadamard's inequality applied to columns, and then to every
cofactor, gives
\[
 1\le d\le3^n,\qquad
 |(C^{-1})_{ij}|\le 3^{n-1}/d.
 \tag{P-F42f}
\]
There are at most $\binom mn\le m^n$ bases.
We choose a single real vector
$\ell=(\ell_1,\ldots,\ell_m)\in[-1,1]^m$.
For each basis and each $i\in I$, put
\[
 w_{J,i}=\ell_i-\ell_J C^{-1}\alpha_i,\qquad
 Z_J=\prod_{i\in I}\max\{1,|w_{J,i}|^{-1}\}.
 \tag{P-F42g}
\]
There exists such a vector for which every $w_{J,i}$ is nonzero
and
\[
                 Z_J\le4m^{2n}9^a
                         \quad\hbox{for every basis }J.
 \tag{P-F42h}
\]
To prove this existence assertion, take the coordinates of
$\ell$ independently and uniformly in $[-1,1]$.
For a uniform random variable $X$ on this interval and any real
number $v$,
\[
 \mathbb E\bigl[\max\{1,|X-v|^{-1}\}^{1/2}\bigr]
 \le1+\frac12\int_{-1}^{1}|u|^{-1/2}\,du=3.
 \tag{P-F42i}
\]
Conditional on $\ell_J$, the variables $w_{J,i}$ for $i\in I$
are independent translates of uniform variables. Hence
$\mathbb E[Z_J^{1/2}]\le3^a$.
Markov's inequality bounds the probability that (P-F42h) fails
for a given basis by $1/(2m^n)$. The union bound over all bases
is at most $1/2$. The probability that any $w_{J,i}$ is zero is
zero, so a simultaneous choice exists. This argument is used
only to prove a bound; no randomized algorithm is asserted.

For this chosen direction, and a shift $b_0=b_\sigma$, set
\[
 \lambda=c\ell_J C^{-1}\mathbf1,\qquad
 v=\ell_J C^{-1}b_0.
\]
The cofactor bounds, the column sum three, and $c\le3$ imply
\[
 |w_{J,i}|\le\frac{(n+1)3^n}{d},\qquad
 |\lambda|\le\frac{n^2 3^n}{d},\qquad
 |v|\le\frac{n^3 3^n}{d}.
 \tag{P-F42j}
\]
For the first inequality use
$1+n3^n/d\le(n+1)3^n/d$, since $d\le3^n$.
Keeping this factor $d$ in the denominator will be useful.

\emph{4. Exact finite-character formula for each tangent cone.}
Let $G_J=\mathbb Z^n/C\mathbb Z^n$, a finite abelian group of
order $d$, and let $\widehat G_J$ be its character group.
At a vertex with basis $J$, the tangent cone retains precisely
the inequalities $\beta_i\ge-\epsilon_i$ for $i\in I$.
Its integer points therefore have $\beta_i\ge0$ for $i\in I$,
and
\[
 \beta_J=C^{-1}\left(cs\mathbf1+b_0-
                            \sum_{i\in I}\alpha_i\beta_i\right).
 \tag{P-F42k}
\]
The remaining coordinates are integers exactly when the vector
in parentheses belongs to $C\mathbb Z^n$.
Finite character orthogonality imposes this congruence.
Consequently the exponential generating function of the cone,
under the specialization $\beta\mapsto e^{z\ell\cdot\beta}$, is
the following meromorphic function near $z=0$:
\[
 \frac1d\sum_{\chi\in\widehat G_J}
 \chi(b_0)^{-1}\chi(c\mathbf1)^{-s}
 \frac{\exp(z(\lambda s+v))}
      {\prod_{i\in I}(1-\chi(\alpha_i)e^{z w_{J,i}})}.
 \tag{P-F42l}
\]
One may first sum the geometric series in a multivariable
domain of convergence for this pointed cone, apply the finite
character filter, and then specialize the resulting rational
identity. Equivalently, before specialization use independent
exponential variables and the rational expressions for
$\beta_J$ in (P-F42k). The formula then follows by analytic
continuation near a generic line. We do not require all
$w_{J,i}$ to have the same sign or claim a common univariate
domain of absolute convergence for all cones.

Brion's identity says that the sum of these rational cone
generating functions over the feasible vertices is the
generating function of the polytope. The rational-polyhedron
version used here is Barvinok \cite{CBarvinokGF}, \emph{The Complexity of
Generating Functions for Integer Points in Polyhedra and
Beyond}, Theorem~2.5, printed p.~9. In particular, the
coefficient of $z^0$ in that sum is $p(cs\mathbf1+b_0)$.
The same identity applies to these lower-dimensional
polyhedra in their ambient lattice; alternatively one can
restrict to their rational affine hull before applying it.

\emph{5. A uniform analytic bound on every character polynomial.}
Set
\[
                      R=\frac1{16n^3 3^n}.
 \tag{P-F42m}
\]
Fix a basis and a character, and let $k$ be the number of
indices $i\in I$ for which $\chi(\alpha_i)=1$.
There are $a$ denominator factors in total.
For $|z|\le R$, (P-F42j) gives
\[
 |zw_{J,i}|\le\frac1{8d},\qquad
 R(|v|+|\lambda|)\le\frac18.
 \tag{P-F42n}
\]
If $\chi(\alpha_i)\ne1$, it is a root of unity of order at
most $d$, so
$|1-\chi(\alpha_i)|\ge4/d$.
The elementary estimate $|e^u-1|\le2|u|$ for $|u|\le1/2$
then gives, throughout this disk,
\[
 \left|\frac1{1-\chi(\alpha_i)e^{z w_{J,i}}}\right|
          \le d/2.
 \tag{P-F42o}
\]
There are no such factors when $d=1$.
For the factors with $\chi(\alpha_i)=1$, the removable
function $z/(1-e^{z w_{J,i}})$ is analytic on this disk and
satisfies
\[
 \left|\frac{z}{1-e^{z w_{J,i}}}\right|
            \le2/|w_{J,i}|.
 \tag{P-F42p}
\]
For example, the latter inequality follows by bounding the
Taylor remainder in $e^u-1=u+O(u^2)$ when $|u|\le1/2$.

Remove the root-of-unity factor $\chi(c\mathbf1)^{-s}$ from
one summand of (P-F42l), and write the rest as
$d^{-1}\chi(b_0)^{-1}z^{-k}e^{\lambda sz}F(z)$, where
\[
 F(z)=e^{vz}
 \prod_{\chi(\alpha_i)=1}
             \frac{z}{1-e^{z w_{J,i}}}
 \prod_{\chi(\alpha_i)\ne1}
             \frac1{1-\chi(\alpha_i)e^{z w_{J,i}}}.
\]
This is analytic on $|z|\le R$.
The constant term of the character summand is
$\chi(c\mathbf1)^{-s}P_{J,\chi,b_0}(s)$, where
$P_{J,\chi,b_0}$ is a polynomial of degree at most $k$.
Writing $\|P\|_1$ for the sum of the absolute values of
the ordinary monomial coefficients, Cauchy's estimate gives
\begin{align*}
 \|P_{J,\chi,b_0}\|_1
 &\le \frac1d
       \sum_{j=0}^k\frac{|\lambda|^j}{j!}
          R^{-(k-j)}\max_{|z|=R}|F(z)|\\
 &\le \frac{2Z_J}{d}\,2^a d^{a-k}R^{-k}\\
 &\le \frac{2Z_J}{d}\,(2/R)^a.
 \tag{P-F42q}
\end{align*}
For the second line, the product of the reciprocal
$|w_{J,i}|$ over any subset is bounded by $Z_J$,
(P-F42o)--(P-F42p) give
$2^k(d/2)^{a-k}\le2^a d^{a-k}$, and
$e^{R(|v|+|\lambda|)}\le2$.
For the last line use $dR\le1$, a consequence of
$d\le3^n$ and (P-F42m).
This absorbs the nontrivial-character denominators into
the same radius bound; multiplying by an additional
factor $d^a$ is unnecessary.

\emph{6. Averaging constituents and completing the estimate.}
For every sufficiently large integer $s$, the same finite
set of feasible bases is used. Formula (P-F42l) and Brion's
identity therefore express $p(cs\mathbf1+b_0)$ as a finite
sum of terms
\[
              \chi(c\mathbf1)^{-s}P_{J,\chi,b_0}(s).
 \tag{P-F42r}
\]
Take a common period for all these roots of unity.
On each residue class, agreement for all sufficiently large
integers identifies the entire eventual constituent
polynomial. Its average over the residue classes retains
exactly the terms with $\chi(c\mathbf1)=1$, by finite
root-of-unity orthogonality. Thus we are averaging
constituent polynomials, not values at a selected set of
integers.
The average polynomial contributed by one basis has norm
at most $2Z_J(2/R)^a$, since there are at most $d$ characters
and each bound in (P-F42q) includes the factor $1/d$.
There are at most $m^n$ bases and $n!$ Weyl shifts.
The triangle inequality in (P-F42c), together with
(P-F42h), now gives
\[
 \begin{split}
 U_n
 &\le n!\,m^n\cdot
          2\bigl(4m^{2n}9^a\bigr)(2/R)^a\\
 &=8n!\,m^{3n}\bigl(288n^3 3^n\bigr)^a.
 \end{split}
 \tag{P-F42s}
\]
The actual function $f_n$ is globally quasipolynomial by
P-F31. Hence the average of the eventual constituents
just computed is exactly its intrinsic polynomial $B_n$
from P-F34: equality on infinitely many integers of
each residue class identifies those polynomials.
Possible cancellations lowering the degree from at most
$a$ to $\binom n3$ are harmless for this upper bound.

Finally, $a=n^3/6+O(n^2)$ and
$\log_2(n!m^{3n})=O(n\log n)$.
Taking logarithms of (P-F42s) proves (P-F42b).
All constants in the proof are uniform in the Weyl shift,
the basis, the finite character, and the stretch parameter.
\end{proof}

\paragraph{Exact remaining quantitative gate.}
P-F34 proves
$\log_2|A_n(e^{2\pi i/2^{n-1}})|=n^4/24+O(n^3)$.
The unconditional coefficient-norm bound above has leading
constant $(\log_2 3)/6>1/24$.
It therefore places the actual norm on an explicit
$2^{O(n^4)}$ scale but leaves the required comparison open.
The proof discards all Weyl signs at its final norm estimate.
Any improvement based on cancellations among the full cubic
weight list and the Weyl numerator requires an additional
argument; neither Brion's identity nor the generic
small-entry bounds alone provide it.

\paragraph{A coefficient-sensitive refinement of the remaining gate.}
Eventual feasibility supplies a stronger estimate for the coefficient
of the stretch parameter than the cofactor estimate used above.
Write \(B_n(X)=\sum_{j=0}^{D}b_{n,j}X^j\), where
\(D=\binom n3\). Then the same construction proves, for every
\(0\le j\le D\),
\[
 |b_{n,j}|
 \le 8n!\,m^{3n}\,18^a
       (16n^3\,3^n)^{a-j}\frac{n^j}{j!},
 \qquad a=m-n.
 \tag{P-F42t}
\]
This estimate concerns the actual coefficients after Weyl
alternation; it does not assume their signs.

Here is the additional proof. A basis that is feasible for all
sufficiently large \(s\) has
\(C^{-1}(c\mathbf1)\ge0\) entrywise. Indeed its basic
coordinates are affine functions of \(s\), with this vector
of slopes, so a negative slope would make that coordinate
negative for all sufficiently large \(s\).
The column sums of \(C\) are three, whence
\(\mathbf1^{\mathsf T}C^{-1}=\mathbf1^{\mathsf T}/3\).
Consequently
\[
 \sum_i (C^{-1}\mathbf1)_i=n/3,
 \qquad
 |\lambda|=|c\ell_JC^{-1}\mathbf1|
              \le cn/3\le n.
 \tag{P-F42u}
\]
This also applies when a limiting basic slope is zero; strict
positivity of the slope is not required.

Use the same radius \(R=(16n^3\,3^n)^{-1}\), direction,
and analytic function \(F\) as in steps~3--5.
For one character summand with \(k\) fixed denominator
factors, the coefficient of \(s^j\) is zero when \(j>k\).
When \(0\le j\le k\), its absolute value is at most
\[
 \frac1d\,\frac{n^j}{j!}\,
       R^{-(k-j)}
       \max_{|z|=R}|F(z)|
 \le
 \frac{2Z_J}{d}\,2^a d^{a-k}
       R^{-(k-j)}\frac{n^j}{j!}
 \le
 \frac{2Z_J}{d}\,2^a
       R^{-(a-j)}\frac{n^j}{j!}.
 \tag{P-F42v}
\]
For the middle inequality, the bound for \(F\) used in
(P-F42q) remains valid, with \(e^{R|v|}\le2\);
no sum over \(j\) is taken here. The last inequality is
exactly \((dR)^{a-k}\le1\).
Averaging retains a subset of at most \(d\) characters.
Summing over at most \(m^n\) bases and \(n!\) Weyl
shifts, and substituting \(Z_J\le4m^{2n}9^a\), gives
(P-F42t). Every basis in that sum is eventually feasible,
so (P-F42u) holds for each of them.

This proves that a definite upper-degree portion of the
average polynomial already lies below the Fourier-amplitude
threshold. More precisely, fix a real constant
\[
 1-\frac1{4\log_2 3}<\theta<1,
 \qquad J_n=\lceil\theta D\rceil,
 \qquad V_{n,\theta}=\sum_{0\le j<J_n}|b_{n,j}|.
 \tag{P-F42w}
\]
Then
\[
 \limsup_{n\to\infty}
 \frac{\log_2\left(\sum_{j=J_n}^{D}|b_{n,j}|\right)}{n^4}
 \le \frac{(1-\theta)\log_2 3}{6}<\frac1{24}.
 \tag{P-F42x}
\]
We use \(\log_2 0=-\infty\) if the indicated tail vanishes.
To verify the estimate, the ratio of the right-hand side
of (P-F42t) at \(j+1\) to that at \(j\) is
\(nR/(j+1)<1\).
The whole tail is therefore at most \(D+1\) times
the bound at \(J_n\). Taking logarithms,
all terms except \(n(a-J_n)\log_2 3\) are
\(O(n^3\log n)\) from above, uniformly in this fixed
\(\theta\); the negative term \(-\log_2(J_n!)\)
can be omitted for this upper estimate.
Since \(a=n^3/6+O(n^2)\) and
\(J_n=\theta n^3/6+O(n^2)\), (P-F42x) follows.

It follows, with the same strict threshold and with no
assumption on a sign or a cancellation, that
\[
 \limsup_{n\to\infty}\frac{\log_2 U_n}{n^4}<\frac1{24}
 \quad\Longleftrightarrow\quad
 \limsup_{n\to\infty}
        \frac{\log_2 V_{n,\theta}}{n^4}<\frac1{24}.
 \tag{P-F42y}
\]
The forward implication uses \(V_{n,\theta}\le U_n\).
For the reverse implication use
\(U_n=V_{n,\theta}+\sum_{j=J_n}^D|b_{n,j}|\),
(P-F42x), and
\(\log_2(x+y)\le1+\max\{\log_2x,\log_2y\}\)
for nonnegative \(x,y\) not both zero.
For example \(\theta=7/8\) is valid: the highest eighth
of the coefficient indices has upper leading constant
\((\log_2 3)/48<1/24\). The unresolved quantitative
gate is thereby confined to the complementary lower
coefficients. Their bound is not supplied here, so the
strict positivity and saturation hypotheses remain open.

\paragraph{Removing pure-cube factors improves the constant and the tail.}
The full determinant of a basis overestimates the orders of
its finite characters when it contains pure-cube columns.
Taking account of this gives the stronger bounds
\[
 \begin{split}
 U_n&\le8n!\,m^{3n}
        \bigl(1152n^4\,5^{n/2}\bigr)^a,\\
 |b_{n,j}|&\le8n!\,m^{3n}\,18^a
        \bigl(64n^4\,5^{n/2}\bigr)^{a-j}
                       \frac{n^j}{j!}
                 \quad(0\le j\le D).
 \end{split}
 \tag{P-F42z}
\]
In particular
\[
 \limsup_{n\to\infty}\frac{\log_2 U_n}{n^4}
             \le\frac{\log_2 5}{12}.
 \tag{P-F42aa}
\]
This is still greater than \(1/24\).

We give all of the changes to the proof, including the distinction
between the order of a finite group and an integer that annihilates
it. Let \(q\) of the selected basis columns be pure cubes.
They are distinct columns \(3e_i\). After permuting rows and
columns, the basis has the block form
\[
 C=\begin{pmatrix}3I_q&B\\0&D_0\end{pmatrix},
 \qquad r=n-q,\qquad d'=|\det D_0|.
 \tag{P-F42ab}
\]
If \(r=0\), use the empty determinant \(d'=1\).
Otherwise \(D_0\) is a nonsingular integer \(r\)-by-\(r\)
matrix. Each of its columns is the restriction of a cubic
column that is not a pure cube, hence has Euclidean norm
at most \(\sqrt5\). Put \(T=5^{n/2}\).
Hadamard's inequality gives \(1\le d'\le5^{r/2}\le T\),
and every cofactor of \(D_0\) has absolute value at most
\(5^{(r-1)/2}\le T\). The latter assertion includes \(r=1\),
whose cofactor is the empty determinant one.
The entries of \(B\) are at most two, so every row sum of
\(B\) is at most \(2r\).
The block inverse
\[
 C^{-1}=
 \begin{pmatrix}
  I_q/3&-BD_0^{-1}/3\\
  0&D_0^{-1}
 \end{pmatrix}
\]
therefore satisfies
\[
       \|C^{-1}\|_{\max}\le \frac{nT}{d'}.
 \tag{P-F42ac}
\]
For the off-diagonal block use
\((2r/3)T/d'\le nT/d'\);
the diagonal \(1/3\) entries obey the same bound because
\(d'\le T\). If \(r=0\), (P-F42ac) follows directly from
\(C^{-1}=I_n/3\).

Set \(E=3d'\). The displayed block inverse shows that
\(EC^{-1}\) is an integer matrix: the off-diagonal block
is \(-B(d'D_0^{-1})\), and the lower block is
\(3d'D_0^{-1}\). Thus \(E\mathbb Z^n\subseteq C\mathbb Z^n\),
so \(E\) annihilates \(G_J=\mathbb Z^n/C\mathbb Z^n\).
Every value of every character of \(G_J\) is consequently
an \(E\)th root of unity.
The number of characters is still
\(d=|\det C|=3^q d'\); it must not be replaced by \(E\)
in the finite filter. The filter retains its factor \(1/d\)
and its sum over \(d\) characters.

Keep the same common generic direction and the same
\(Z_J\) bound. From (P-F42ac), the column sum three,
and \(|(b_0)_i|\le n-1\), we obtain
\[
 |w_{J,i}|\le\frac{4n^2T}{d'},\qquad
 |v|\le\frac{n^4T}{d'},\qquad |\lambda|\le n.
 \tag{P-F42ad}
\]
For the first inequality,
\(|\ell_JC^{-1}\alpha_i|\le3n\|C^{-1}\|_{\max}\);
the additional one is at most \(T/d'\).
For the second, use
\(\sum_{i,j}|\ell_i(C^{-1})_{ij}(b_0)_j|
 \le n^2(n-1)\|C^{-1}\|_{\max}\).
The last inequality is the feasible-slope result (P-F42u).
Take
\[
                  R'=(64n^4T)^{-1}.
 \tag{P-F42ae}
\]
For \(n\ge4\), these estimates give
\[
 R'|w_{J,i}|\le\frac{3}{16n^2E}\le\frac1{8E},
 \qquad R'(|v|+|\lambda|)\le\frac1{32},
 \qquad ER'\le\frac{3}{64n^4}\le1.
 \tag{P-F42af}
\]
Every nontrivial character value \(\alpha\) satisfies
\(|1-\alpha|\ge4/E\).
The analytic proof in step~5 therefore applies on
\(|z|\le R'\) with the nontrivial denominator bound \(E/2\),
while a fixed denominator still has bound
\(2/|w_{J,i}|\) after multiplication by \(z\).
For the coefficient of \(s^j\), the same Cauchy calculation
now reads
\[
 \bigl|[s^j]P_{J,\chi,b_0}\bigr|
 \le
 \frac{2Z_J}{d}\,2^a E^{a-k}
           (R')^{-(k-j)}\frac{n^j}{j!}
 \le
 \frac{2Z_J}{d}\,2^a
           (R')^{-(a-j)}\frac{n^j}{j!},
 \tag{P-F42ag}
\]
with coefficient zero for \(j>k\).
The last inequality uses \((ER')^{a-k}\le1\).
When bounding the full coefficient norm before summing,
the Taylor coefficients of \(e^{\lambda sz}\) contribute
at most \(e^{R'|\lambda|}\); together with the numerator
factor \(e^{R'|v|}\) this is at most two by (P-F42af).
Hence the full norm per character is at most
\((2Z_J/d)(2/R')^a\), just as in the original proof.
There are at most \(d\) retained characters, at most \(m^n\)
bases, and \(n!\) shifts. Substituting
\(Z_J\le4m^{2n}9^a\) proves both estimates in (P-F42z),
with \(18/R'=1152n^4T\).
Taking logarithms proves (P-F42aa).

The tail and equivalence arguments (P-F42w)--(P-F42y)
now hold with the weaker requirement
\[
       1-\frac1{2\log_2 5}<\theta<1,
 \qquad
 \limsup_{n\to\infty}
 \frac{\log_2\left(\sum_{j=\lceil\theta D\rceil}^{D}
                            |b_{n,j}|\right)}{n^4}
       \le\frac{(1-\theta)\log_2 5}{12}<\frac1{24}.
 \tag{P-F42ah}
\]
Indeed the ratio of the coefficient bounds at successive
indices is \(nR'/(j+1)<1\), and the only term of order
\(n^4\) in their logarithms is
\((a-j)(n/2)\log_2 5\).
For the simple choice \(\theta=4/5\), the highest fifth
has upper leading constant \((\log_2 5)/60<1/24\).
Thus the strict \(1/24\) bound for the actual full norm
\(U_n\) is equivalent to that bound for
\(\sum_{j<\lceil4D/5\rceil}|b_{n,j}|\).
This is a smaller unresolved set of coefficients;
no bound of the required size is proved for it here.

\paragraph{The corresponding signed test at polynomial arguments.}
The controlled tail can also be removed from the exact
pointwise inequality in P-F34, without replacing the
remaining polynomial by its coefficient norm.
Put
\[
 J_n=\lceil4D/5\rceil,\qquad
 L_n(X)=\sum_{j<J_n}b_{n,j}X^j,\qquad
 \eta=\frac1{24}-\frac{\log_2 5}{60}>0.
 \tag{P-F42ai}
\]
For every fixed real \(K\ge0\), all sufficiently large \(n\)
satisfy the uniform estimate
\[
 \sup_{0\le x\le n^K}|B_n(x)-L_n(x)|
       \le 2^{-\eta n^4/2}\,
                 |A_n(e^{2\pi i/2^{n-1}})|.
 \tag{P-F42aj}
\]
Indeed the left side is at most
\(n^{KD}\sum_{j=J_n}^D|b_{n,j}|\).
The proof of (P-F42ah) bounds its base-two logarithm by
\((\log_2 5)n^4/60+O_K(n^3\log n)\).
The amplitude formula (P-F34t) has logarithm
\(n^4/24+O(n^3)\), so their ratio has logarithm at most
\(-\eta n^4+O_K(n^3\log n)\), which is at most
\(-\eta n^4/2\) for all sufficiently large \(n\).

In particular, if either \(x=p(f_n)\) or \(x=s(f_n)+1\)
is at most \(n^K\), its necessary Fourier inequality gives
\[
 L_n(x)\ge
       (1-2^{-\eta n^4/2})
                 |A_n(e^{2\pi i/2^{n-1}})|.
 \tag{P-F42ak}
\]
Thus a bound of the form
\(\sup_{0\le x\le n^K}L_n(x)\le(1-\epsilon_K)|A_n|\),
with some fixed \(\epsilon_K>0\) for every fixed \(K\)
and all sufficiently large \(n\), would also exclude
polynomially bounded strict indices. This is a signed
value comparison and can be useful even if the coefficient
norm has limiting quartic exponent \(1/24\).
For example it would suffice to have
\(b_{n,0}\le(1-2\epsilon_K)|A_n|\) and
\(\sum_{1\le j<J_n}|b_{n,j}|n^{Kj}\le\epsilon_K|A_n|\)
eventually for each fixed \(K\).
These remain sufficient tests for excluding small indices;
the necessary Fourier inequality is not an equivalence
characterizing either index.
No such comparison, or corresponding upper bound on
the constant term \(b_{n,0}\), is established here.
